Reactive autoscaling is easy to deploy but often responds too late. This study evaluates a reproducible predictive alternative on public Alibaba traces using aggregate, machine-level, transfer, and broad service-level evidence. The pipeline couples leakage-free multi-horizon forecasting with a trace-driven controller that consumes median and upper-quantile predictions. On the aggregate series, linear regression achieves the lowest average MAE (3.514), while \texttt{UA-MSTCN-Lite} remains competitive and provides calibrated $P90$ coverage of 0.920, 0.867, and 0.868 at the 1-, 5-, and 10-minute horizons. A 12-machine cohort and a 3-fold zero-shot cross-machine benchmark confirm that simple linear structure remains difficult to beat for point prediction. A broad audited service cohort of 139 \texttt{app\_du} groups, extracted from the official container tables and a full mirrored pass over \texttt{container\_usage}, changes the service-level headline: linear regression becomes the strongest average point forecaster (MAE 32.571), while \texttt{UA-MSTCN-Lite} retains useful $P90$ coverage of 0.921, 0.896, and 0.877. Against reactive control, predictive control raises mean SLA violation from 0.017 to 0.031 and scaling actions from 249.0 to 399.4; mean over-provisioning falls from 56.31 to 53.68, but the paired service bootstrap confidence interval crosses zero. Public traces therefore support reproducible predictive-autoscaling evaluation and boundary mapping, but not a claim that the current shared global predictive policy is uniformly better across services.
